% TOP_PROBLEM: {"id": 4700020, "title": "A Markus–Yamabe/La Salle problem for discrete dynamical systems", "category_id": 11, "category": {"id": 11, "name": "dynamical_systems", "display_name": "Dynamical Systems", "description": "Problems about long-term behavior of deterministic systems, Hamiltonian dynamics, and periodic orbits.", "slug": "dynamical-systems", "order_index": 11, "created_at": "2026-07-31T15:26:25.671Z"}, "status": "open", "problem_number": "AMR-046-0020", "set_id": 13}
% FIRST_POSTED: 2026-10-02
% ENTRY_KIND: statement_only
% UnsolvedMath ID 4700020; source catalogue code AMR-046-0020
% !TeX program = lualatex
% Definition and sources only; this file is not an LLM solution attempt.
\documentclass[11pt,a4paper]{article}
\usepackage[margin=25mm]{geometry}
\usepackage{fontspec}
\setmainfont{lmroman10-regular.otf}[BoldFont=lmroman10-bold.otf,
 ItalicFont=lmroman10-italic.otf,BoldItalicFont=lmroman10-bolditalic.otf]
\usepackage{amsmath,amssymb,mathtools}
\usepackage{unicode-math}
\setmathfont{latinmodern-math.otf}
\AtBeginDocument{\let\setminus\smallsetminus}
\usepackage{xurl,hyperref}
\hypersetup{colorlinks=true,urlcolor=blue}
\setcounter{secnumdepth}{0}
\setlength{\parindent}{0pt}
\setlength{\parskip}{5pt}
\setlength{\emergencystretch}{3em}
\begin{document}
\section*{A Markus–Yamabe/La Salle problem for discrete dynamical systems}\label{4700020}
{\small UnsolvedMath ID 4700020 · AMR-046-0020 · Dynamical Systems · AMR Open Problem Lists}
\subsection{Definitions and mathematical statement}
Let $F:\mathbb R^2\to\mathbb R^2$ be smooth and suppose $F(a)=a$ for some
$a\in\mathbb R^2$. For a real matrix $M=(m_{ij})$, write $|M|=(|m_{ij}|)$
for entrywise absolute value, and let $\rho(M)$ be its spectral radius,
the greatest modulus of its complex eigenvalues. Assume
\[
 \rho\bigl(|DF(x)|\bigr)<1\qquad\text{for every }x\in\mathbb R^2.
\]
The question is whether $a$ must be globally asymptotically stable for the
discrete dynamical system $x_{j+1}=F(x_j)$.

Here global asymptotic stability requires both
\[
 \lim_{j\to\infty}F^j(x)=a\quad\text{for all }x\in\mathbb R^2
\]
and Lyapunov stability: for every $\varepsilon>0$ there is $\delta>0$ such
that $\|x-a\|<\delta$ implies $\|F^j(x)-a\|<\varepsilon$ for all
integers $j\ge0$. All norms on $\mathbb R^2$ give the same qualitative
definition.

The absolute value in the hypothesis is essential: it is taken before
computing the spectral radius, not afterwards. Nor is the condition a
uniform contraction estimate in a single chosen norm. The permitted
spectral radii may approach one as $x$ goes to infinity, and no global
bound on $DF$ is assumed. Thus a proof that invokes a fixed contraction
constant must first derive it, rather than silently adding it as a
hypothesis. The fixed point is assumed to exist; its uniqueness would be
part of the asserted conclusion.
\subsection{Short English statement}
Does pointwise spectral-radius control of the entrywise absolute Jacobian force global convergence to the fixed point?
\subsection{Sources}
\begin{itemize}
\item[S1] ULAM.AI and the UnsolvedMath Contributors, supplied problems.json, record 4700020 (AMR-046-0020), AMR Open Problem Lists. Catalogue identity and source transcription; CC BY 4.0.
\url{https://huggingface.co/datasets/ulamai/UnsolvedMath}.
\item[S2] Gasull - Some open problems in low dimensional dynamical systems (2020), Problem environment 20: A Markus–Yamabe/La Salle problem for discrete dynamical systems. Original source indexed by AMR-046-0020.
\url{https://arxiv.org/abs/2012.02524}.
\end{itemize}
\end{document}
